\section{Limitations}
\label{sec:limitations}

\textbf{Per-head trajectories: single seed.}
The per-head L1H6 trajectories (Figures~5--7, the IS values in
Table~\ref{tab:main-results}) are from seed 42 only. Seeds 123 and 7 were
trained (checkpoints exist for both conditions) and their \emph{aggregate}
metrics (mean induction score over all 16 heads, training loss) were
successfully extracted and are reported in Table~\ref{tab:main-results} as
3-seed mean $\pm$ SD. However, the per-head checkpoint sweep
(\texttt{sweep\_checkpoints()}, which yields the full $[\text{step},
\text{layer}, \text{head}]$ tensor needed for L1H6-specific values) was run
for seed 42 only. Running this sweep for seeds 123 and 7 is the
single highest-priority remaining step to obtain per-head confidence
intervals and statistically confirm the reinforcement direction.

\textbf{Three checkpoints per run.}
With \texttt{checkpoint\_every=100} and 244 total steps, each run produces
checkpoints at steps 0, 100, 200 only. The Savitzky-Golay smoothing window
(window $= 5$) exceeds this series length and falls back to raw consecutive
differences, limiting the granularity of phase-transition detection. More
frequent checkpointing would allow finer-grained trajectory analysis.

\textbf{Single-head circuit.}
Only L1H6 exceeds the $\geq 0.5$ attribution threshold; no layer-0
previous-token partner was identified. The canonical induction circuit
in \citet{olsson2022context} consists of a previous-token head (layer 0)
$+$ induction head (layer 1) pair. Our activation patching identifies the
layer-1 head as the dominant causal component, but the absence of a
confirmed layer-0 partner means the circuit topology is less complete than
described in prior work. This may reflect the specific threshold choice or
the pooled corruption used in patching.

\textbf{Model scope.}
All results are from a 2-layer, attention-only transformer (no MLP layers).
Claims are scoped to this model class; models with MLP layers, residual
stream interventions via both attention and MLP, may show qualitatively
different circuit stability properties.

\textbf{Dataset change.}
Between experimental sessions, the primary code dataset was changed from
\texttt{codeparrot/github-code} to \texttt{transformersbook/codeparrot}
due to a breaking change in the \texttt{datasets} library (DECISION-013).
Both datasets derive from the same GitHub BigQuery source, but the subset
filtering differs slightly. Strict reproducibility requires running
notebook 03 with the v9 codebase.

\textbf{Circuit threshold.}
The attribution threshold of $0.5$ (DECISION-002, stated a priori) is
arbitrary. The IS threshold of $0.5$ shown in Figures~5 and~6 is a
visualisation reference, not a hard criterion. Raw values are reported
throughout to allow alternative thresholds.

\textbf{Token budget.}
500K tokens may not capture changes emerging at larger budgets. The apparent
plateau between steps 100 and 200 ($\Delta = 0.054$ for code, $0.033$ for
prose) raises the question of whether longer runs would continue, reverse,
or accelerate the IS increase.

\textbf{Adversarial probe normalisation.}
The \texttt{all\_same\_token} probe produces a normalised logit difference of
$4.12$ (pre) and $2.41$ (post), both exceeding the clean baseline of $1.0$.
This reflects a normalisation artefact: when every token is identical, every
position trivially satisfies the prefix-matching criterion, making the
numerator abnormally large. This result should not be interpreted as evidence
of ``super-induction''; it is an artefact of the degenerate input
distribution interacting with the normalisation scheme.
